% =============================================================================
% CORTEX: A Unified Cognitive Tree Architecture for Persistent, Navigable,
% and Self-Consolidating LLM Agent Memory
%
% arXiv-style preprint. Compile with pdflatex (twice) or latexmk -pdf.
% =============================================================================
\documentclass[11pt,a4paper]{article}

% ---- Page geometry & typography --------------------------------------------
\usepackage[margin=1in]{geometry}
\usepackage[utf8]{inputenc}
\usepackage[T1]{fontenc}
\usepackage{lmodern}
\usepackage{microtype}
\usepackage{parskip}

% ---- Math, theorems, symbols -----------------------------------------------
\usepackage{amsmath,amssymb,amsthm}

% ---- Tables ----------------------------------------------------------------
\usepackage{booktabs}
\usepackage{array}
\usepackage{tabularx}
\usepackage{multirow}
\usepackage{longtable}

% ---- Graphics, color, hyperlinks -------------------------------------------
\usepackage{graphicx}
\usepackage[dvipsnames]{xcolor}
\usepackage[
  colorlinks=true,
  linkcolor=blue!50!black,
  citecolor=blue!50!black,
  urlcolor=blue!60!black,
  pdftitle={CORTEX: A Unified Cognitive Tree Architecture},
  pdfauthor={Anonymous}
]{hyperref}
\usepackage{url}

% ---- Lists, captions, headers ----------------------------------------------
\usepackage{enumitem}
\setlist[itemize]{leftmargin=*,topsep=2pt,itemsep=2pt}
\setlist[enumerate]{leftmargin=*,topsep=2pt,itemsep=2pt}
\usepackage{caption}
\captionsetup{font=small,labelfont=bf}

% ---- Code listings ---------------------------------------------------------
\usepackage{listings}
\definecolor{codebg}{rgb}{0.97,0.97,0.97}
\definecolor{codecmt}{rgb}{0.4,0.4,0.4}
\definecolor{codekey}{rgb}{0.0,0.0,0.6}
\definecolor{codestr}{rgb}{0.0,0.5,0.0}
\lstset{
  basicstyle=\ttfamily\footnotesize,
  backgroundcolor=\color{codebg},
  commentstyle=\color{codecmt}\itshape,
  keywordstyle=\color{codekey}\bfseries,
  stringstyle=\color{codestr},
  showstringspaces=false,
  breaklines=true,
  frame=single,
  framerule=0pt,
  framesep=4pt,
  xleftmargin=4pt,
  xrightmargin=4pt,
  aboveskip=8pt,
  belowskip=8pt,
}
\lstdefinelanguage{Python}{
  keywords={async, await, def, return, if, else, elif, for, while, in, not,
    and, or, is, None, True, False, class, import, from, try, except,
    raise, with, as, lambda, yield, assert, pass, break, continue, global,
    nonlocal, self},
  morecomment=[l]{\#},
  morestring=[b]",
  morestring=[b]',
}
\lstdefinelanguage{SQL}{
  keywords={SELECT, FROM, WHERE, AND, OR, JOIN, LEFT, RIGHT, ON, INSERT,
    UPDATE, DELETE, INTO, VALUES, AS, ORDER, BY, GROUP, HAVING, LIMIT,
    UNION, ALL, WITH, RECURSIVE, RETURNING, NULL, IS, NOT, IN, SET,
    INTERVAL, NOW, COALESCE, GREATEST, ANY, CAST, EXISTS, CASE, WHEN,
    THEN, END, DISTINCT, MAX, MIN, COUNT, AVG, SUM},
  morecomment=[l]{--},
  morestring=[b]'
}

% ---- Theorem-like environments for formal properties -----------------------
\newtheoremstyle{cortexprop}%
  {8pt}{8pt}%        space above/below
  {}%                body font
  {}%                indent
  {\bfseries}%       head font
  {.}%               head punctuation
  {0.5em}%           space after head
  {\thmname{#1}~\thmnumber{#2}\thmnote{ (\textit{#3})}}
\theoremstyle{cortexprop}
\newtheorem{property}{Property}

% ---- Title block -----------------------------------------------------------
\title{
  \vspace{-1.5em}
  \textbf{CORTEX: A Unified Cognitive Tree Architecture for\\
  Persistent, Navigable, and Self-Consolidating LLM Agent Memory}
  \vspace{0.5em}
}
\author{
  Author One$^{1}$ \quad Author Two$^{1}$ \quad Author Three$^{1}$\\[0.5em]
  \small $^{1}$\textit{Affiliation, Address}\\
  \small Correspondence: \texttt{author1@org.tld}
}
\date{\small Preprint --- May 2026}

% ---- Custom commands -------------------------------------------------------
\newcommand{\cortex}{\textsc{Cortex}}
\newcommand{\op}[1]{\textsc{#1}}
\newcommand{\code}[1]{\texttt{#1}}

% =============================================================================
\begin{document}
\maketitle

% =============================================================================
\begin{abstract}
\noindent
Long-running LLM agents face a structural tension: their reasoning is bounded
by a finite context window, yet useful tasks demand continuity across hours,
days, and many tool invocations. Existing agent memory systems address parts
of the problem --- episodic logs (Mem0, Zep), tiered context management
(MemGPT/Letta), reflection trees (Generative Agents), reasoning-based
retrieval (PageIndex), and recursive task decomposition (RLM) --- but no
public framework integrates a writable cognitive workspace, persistent
multi-domain memory, automatic consolidation, sub-agent isolation, and an
associative graph overlay into a single coherent runtime. We present
\textsc{Cortex} (Cognitive Orchestrated Recursive Tree EXecution), an
agent-memory architecture in which the agent's complete cognitive state is a
typed, persistent tree, and its language-model context window is merely a
bounded \textit{viewport} onto that tree. \textsc{Cortex} unifies four memory
domains --- Knowledge, Episodic, Experience, and Intelligence --- onto a
single ORM model trio (\code{tree / node / edge}), disambiguated by
enumeration. A background ``Dreaming'' pipeline consolidates episodes into
observations, observations into patterns, and patterns into actionable rules,
mirroring established cognitive-science models. A semantic graph layer
overlaying all trees enables hybrid embedding-plus-graph retrieval with
auto-decay maintenance. We describe the data model, the seven primitive
operations, the consolidation pipeline, the multi-tenant scope hierarchy,
integration with an execution engine, and the system's invariants. We
characterize \textsc{Cortex}'s design properties, position it against the
contemporary memory-framework landscape, and discuss limitations and open
questions.
\end{abstract}

\smallskip
\noindent\textbf{Keywords:} large language models, agent memory, cognitive
architecture, hierarchical retrieval, knowledge consolidation, multi-agent
systems, context engineering.

% =============================================================================
\section{Introduction}
\label{sec:intro}

The capabilities of large language model (LLM) agents have advanced rapidly,
yet four structural problems persist across deployed systems:

\begin{enumerate}
  \item \textbf{Context rot.} Empirical and anecdotal evidence shows that
  model accuracy degrades as prompt length grows, even within stated context
  limits --- a phenomenon Anthropic terms ``context rot.'' Filling the window
  with running transcripts is therefore not a viable long-task strategy.

  \item \textbf{Opaque retrieval.} Vector-based retrieval-augmented generation
  (RAG) flattens document hierarchy into embeddings, sacrificing structural
  information and provenance. Recent work (PageIndex, Vectify AI 2025)
  demonstrates that reasoning-based retrieval over a hierarchical document
  tree can lift FinanceBench accuracy from approximately 31\% to 98.7\%.

  \item \textbf{Cross-run amnesia.} Most production agents start each
  invocation from a blank slate. Persistent memory frameworks like Mem0 and
  Zep address this for facts but do not provide a writable workspace the
  agent can edit and revisit.

  \item \textbf{No learning loop.} Systems that record execution history
  typically lack a mechanism to abstract that history into reusable rules.
  The cognitive-science literature on \textit{memory consolidation} ---
  episodic experience compressed into semantic knowledge --- has obvious
  parallels but is implemented in few production systems.
\end{enumerate}

We argue that addressing these four problems requires not a new model, but a
new \textit{runtime substrate} for the agent's cognitive state. Specifically,
we propose that the agent's working state should be modeled as a persistent,
navigable, writable typed tree, and that the language model itself should
only ever see a bounded slice --- a \textit{viewport} --- of that tree. All
other context is reachable by navigation, paging, or sub-task delegation.

\paragraph{Contributions.} This paper makes the following contributions:

\begin{enumerate}
  \item \textbf{A unified data model} in which four orthogonal memory domains
  (Knowledge, Episodic, Experience, Intelligence) and a transient runtime
  workspace share a single ORM trio (\code{tree / node / edge}),
  disambiguated by enumeration. One set of query, retrieval, graph, and
  embedding primitives serves every domain (\S\ref{sec:overview},
  \S\ref{sec:datamodel}).

  \item \textbf{Seven primitive operations} (\op{Navigate}, \op{Read},
  \op{Write}, \op{Recurse}, \op{Await\_Children}, \op{Checkpoint},
  \op{Assemble}) that an agent invokes to manipulate its cognitive tree,
  formalizing the ``agent-as-tree-explorer'' pattern. We elevate these to
  first-class planner step types rather than hiding them behind tool calls
  (\S\ref{sec:ops}).

  \item \textbf{A Dreaming consolidation pipeline} that runs the cognitive
  sequence \textit{episode $\to$ observation $\to$ pattern $\to$ rule} as a
  typed three-phase background job with confidence and recurrence
  thresholds, embedding-based clustering, and a generation counter
  (\S\ref{sec:dreaming}).

  \item \textbf{A semantic-graph overlay} (\code{cortex\_edges}) enabling
  associative cross-tree, cross-domain retrieval with auto-edge creation on
  embedding, weight decay over inactivity, and co-access tracking during
  execution (\S\ref{sec:graph}).

  \item \textbf{A six-level scope hierarchy} (App $\to$ Partner $\to$ Tenant
  $\to$ User $\to$ Entity $\to$ Runtime) that unifies platform-wide,
  multi-tenant, per-user, per-agent, and per-run state within the same
  schema (\S\ref{sec:scope}).

  \item \textbf{Subtree-isolated recursive children} via a recursive-CTE
  ancestry check, solving the ``sub-agent leaking context to parent''
  failure mode common in multi-agent systems (\S\ref{sec:isolation}).
\end{enumerate}

We position \textsc{Cortex} against contemporary frameworks in
\S\ref{sec:related}, describe the architecture and implementation in
\S\ref{sec:overview}--\S\ref{sec:impl}, characterize design properties in
\S\ref{sec:discussion}, discuss limitations in \S\ref{sec:limits}, and
outline future work in \S\ref{sec:future}.

% =============================================================================
\section{Related Work}
\label{sec:related}

\subsection{Tiered and Operating-System-Inspired Memory}
MemGPT~\cite{packer2023memgpt}, now developed as Letta, treats the LLM as an
operating system with core memory, archival memory, and recall memory tiers,
paging information between tiers via tool calls. \textsc{Cortex} adopts the
paging idea but extends it from a tiered store to a typed tree, adds four
orthogonal domains rather than three abstract tiers, and supplements the
agent-driven paging with system-driven background consolidation.

\subsection{Memory-as-a-Service Frameworks}
Mem0~\cite{bansal2025mem0} extracts facts from conversations into a hybrid
vector + graph + key-value store with user, session, and agent scopes.
Zep introduces a temporal knowledge graph capturing state changes (``I used
to live in London, but I moved to Tokyo''). Cognee emphasizes deep knowledge
graphs. These systems excel at fact recall but do not provide a writable
workspace, recursive sub-task scoping, or output canvas.

\subsection{Reflection and Generative Agents}
Park et al.~\cite{park2023generative} introduced reflection trees in which
agents periodically summarize the latest ${\sim}100$ memories into
higher-level insights, retrieved via a weighted combination of recency,
importance, and relevance. \textsc{Cortex}'s Dreaming Engine generalizes
this from a single LLM reflection prompt to a typed three-phase pipeline
with explicit confidence/strength thresholds and edge creation between
phases.

\subsection{Cognitive-Architecture-Inspired Frameworks}
The Cognitive Architectures for Language Agents
framework~\cite{sumers2023coala} (CoALA) provides the canonical taxonomy:
working, episodic, semantic, and procedural memory. \textsc{Cortex} is a
concrete instantiation of CoALA with one refinement: it splits CoALA's
``semantic + procedural'' into three orthogonal domains (Knowledge /
Experience / Intelligence), with Experience serving as the explicit
consolidation midpoint between raw episodes and actionable rules. The
classical cognitive architectures Soar~\cite{laird2022soar} and ACT-R
provide the broader inspiration --- production rules, declarative memory,
and chunking.

\subsection{Reasoning-Based Retrieval}
PageIndex~\cite{vectifyai2025pageindex} demonstrates that hierarchical tree
navigation by an LLM agent outperforms vector search on structured
documents. \textsc{Cortex} adopts PageIndex's bounded-viewport navigation
primitive but extends the tree from a read-only document index to a
writable workspace that persists across runs and overlays a vector + graph
layer for hybrid retrieval rather than purely-reasoning retrieval.

\subsection{Recursive Language Models}
Recursive Language Models (RLM)~\cite{zhang2025rlm} let the LLM treat the
prompt as a variable and recursively sub-query itself, achieving robust
performance at 10M+ token inputs. \textsc{Cortex}'s \op{Recurse} primitive
applies the same idea to a typed tree: a child run is spawned scoped to a
subtree (\code{scoped\_subtree\_root\_id}) and results are collected via
\op{Await\_Children}. The tree itself is the bounded environment.

\subsection{Hierarchical Multi-Agent Memory}
G-Memory~\cite{wang2025gmemory} introduces a three-tier hierarchical memory
(insight / query / interaction graphs) for multi-agent systems.
ByteRover~\cite{byterover2026} represents knowledge in a Domain $\to$ Topic
$\to$ Subtopic $\to$ Entry context tree with importance scoring and maturity
tiers. \textsc{Cortex} overlaps these on importance scoring
(\code{importance\_score}), generation counters
(\code{consolidation\_generation}), and decay (\code{access\_count}, edge
weight decay), while broadening scope coverage and integrating tightly with
an execution engine.

\subsection{Filesystem-as-Memory}
Cognitive Workspace~\cite{liu2025cogworkspace} and Letta's
filesystem-benchmark work argue that a plain filesystem with markdown notes
can match or beat vector RAG for agents that manage their own scratchpad.
\textsc{Cortex} can be viewed as a \textit{schematized} filesystem: every
``file'' carries type, status, embedding, parent, summary, content,
provenance, importance, and edges. The structure enables both LLM
navigation and SQL queries --- a property a plain filesystem lacks.

\subsection{Summary of Positioning}
Individual ideas in \textsc{Cortex} have published antecedents. The
contribution is an \textbf{integrated runtime}: to our knowledge, no public
framework simultaneously provides (a) a single schema unifying knowledge,
episodic, experience, intelligence, and runtime working state; (b)
six-level scope inheritance; (c) automatic Dreaming-style consolidation;
(d) viewport-based bounded context with paged reads; (e) recursive
subtree-isolated children; and (f) a graph layer with auto-edge creation,
decay, and co-access tracking, behind one ORM model and seven primitive
operations.

% =============================================================================
\section{System Overview}
\label{sec:overview}

\subsection{Core Thesis}

\begin{quote}
\itshape
The agent's complete cognitive state is the tree. The language model's
context window is a viewport onto the tree.
\end{quote}

Every observation, finding, knowledge reference, sub-task, checkpoint, and
output paragraph the agent produces is a node in a persistent typed tree.
The LLM's input at any step is the rendered viewport of the current cursor
node --- its title, summary, immediate children's summaries, parent, and
breadcrumb --- plus an enumeration of the seven \textsc{Cortex} operations
the agent can invoke.

\subsection{Logical Layering}

\begin{lstlisting}[basicstyle=\ttfamily\scriptsize]
+---------------------------------------------------------------------+
|   Execution Engine (DAG / Recursive planner)                        |
|                                                                     |
|   CortexBridge   ---->  CortexRouter (7 ops)  ---->  PostgreSQL +   |
|   (interface)           MemoryAssembler              pgvector       |
|                         DreamingEngine                              |
|                         SemanticGraphService                        |
+---------------------------------------------------------------------+
\end{lstlisting}

\subsection{The Four Memory Domains}

\textsc{Cortex} persists four orthogonal long-lived trees per agent entity,
plus one transient runtime tree per execution
(Table~\ref{tab:domains}).

\begin{table}[h]
\centering
\caption{The four persistent memory domains and the transient runtime
workspace.}
\label{tab:domains}
\small
\begin{tabularx}{\linewidth}{lXXX}
\toprule
\textbf{Domain} & \textbf{Source} & \textbf{Written by} & \textbf{Used for}\\
\midrule
Knowledge & User documents, scraped data & Ingestion pipelines & Reference recall \\
Episodic & Completed execution runs & System (on run completion) & Recent history \\
Experience & Distilled from episodes & Dreaming Engine (phases 1--2) & Observations and patterns \\
Intelligence & Distilled from patterns & Dreaming Engine (phase 3) & Actionable rules \\
\midrule
Runtime \textit{(transient)} & The current execution & Agent and execution engine & Working memory, current output \\
\bottomrule
\end{tabularx}
\end{table}

\subsection{Design Principles}

\textsc{Cortex} is built on four invariants that are enforced in code:

\begin{enumerate}
  \item \textbf{Summary Always Exists.} A node cannot have children unless
  its \code{summary} is set. This guarantees a parent's viewport is always
  informative.
  \item \textbf{No Unbounded Viewports.} A node's direct-children count is
  bounded by \code{max\_children} (default 12). Exceeding the bound triggers
  asynchronous re-clustering into a group node.
  \item \textbf{Content is Always Paged.} Large \code{content} fields are
  read in slices of \code{page\_size\_tokens} (default 8000 tokens). No
  single read returns more than one page.
  \item \textbf{Write-Once Content.} Node content is immutable; revisions
  are modeled as child nodes, not in-place edits. This makes provenance
  traceable and prevents subtle race conditions.
\end{enumerate}

% =============================================================================
\section{Data Model}
\label{sec:datamodel}

\textsc{Cortex} uses three tables stored in PostgreSQL with the
\code{pgvector} extension.

\subsection{The \code{cortex\_trees} Table}

Each tree row is a cognitive context. Key columns:

\begin{itemize}
  \item \code{id}, \code{entity\_id}, \code{user\_id}, \code{company\_id} ---
  ownership and tenancy.
  \item \code{task\_description}, \code{status} $\in$ \{active, suspended,
  complete, archived\}.
  \item \code{root\_node\_id}, \code{output\_root\_id},
  \code{resume\_cursor\_id} --- cursors into the node table.
  \code{resume\_cursor\_id} is the single source of truth for ``where the
  agent is.''
  \item \code{max\_children} (default 12), \code{page\_size\_tokens}
  (default 8000), \code{context\_budget\_pct} (default 40) --- invariants
  and tunables.
  \item \code{memory\_domain} $\in$ \{knowledge, episodic, experience,
  intelligence\}, \code{scope\_level} $\in$ \{app, partner, tenant, user,
  entity, runtime\} --- the two enums that disambiguate every tree.
  \item \code{last\_consolidated\_at}, \code{consolidation\_generation},
  \code{source\_run\_ids} --- Dreaming metadata.
  \item \code{resume\_schedule}, \code{next\_resume\_at} --- multi-day
  scheduled wake-ups.
\end{itemize}

\subsection{The \code{cortex\_nodes} Table}

Every piece of information is a node. Key columns:

\begin{itemize}
  \item \code{id}, \code{tree\_id}, \code{parent\_id} --- structural.
  \item \code{node\_type} --- one of 18 typed values: \code{root},
  \code{knowledge}, \code{finding}, \code{task}, \code{output},
  \code{checkpoint}, \code{group}, \code{document}, \code{section},
  \code{chunk}, \code{observation}, \code{pattern}, \code{suggestion},
  \code{instruction}, \code{strategy}, \code{preference}, \code{episode},
  \code{episode\_group}.
  \item \code{title}, \code{summary}, \code{content},
  \code{content\_tokens} --- display and content.
  \item \code{status} $\in$ \{pending, active, complete, summarised\}.
  \item \code{source\_ref}, \code{execution\_run\_id},
  \code{metadata\_extra} --- provenance.
  \item \code{embedding} (\code{pgvector.Vector(768)}),
  \code{embedding\_model} --- semantic vector.
  \item \code{cross\_refs}, \code{access\_count}, \code{last\_accessed\_at},
  \code{importance\_score} --- graph and learning signals.
\end{itemize}

\subsection{The \code{cortex\_edges} Table}

Directed, weighted, typed edges overlaying all trees. Edges may cross
trees, domains, and scope levels.

\begin{itemize}
  \item \code{source\_node\_id}, \code{target\_node\_id},
  \code{edge\_type}, \code{weight} $\in [0,1]$.
  \item \code{traversal\_count}, \code{last\_traversed\_at} --- for
  boosting and decay.
  \item \code{created\_by} --- provenance string (e.g.,
  \code{"dreaming\_engine"}, \code{"embedding\_pipeline"}).
  \item Uniqueness: \code{(source, target, edge\_type)}. Inserts are
  upserts.
\end{itemize}

The edge-type vocabulary is given in Table~\ref{tab:edgetypes}.

\begin{table}[h]
\centering
\caption{Edge-type vocabulary in the semantic graph layer.}
\label{tab:edgetypes}
\small
\begin{tabularx}{\linewidth}{lXl}
\toprule
\textbf{Edge type} & \textbf{Semantics} & \textbf{Created by}\\
\midrule
\code{semantic\_similar} & Cosine similarity $\geq 0.85$ & Auto on embedding \\
\code{derived\_from} & Pattern derived from observation & Dreaming phase 2 \\
\code{generalizes} & Rule generalizes a pattern & Dreaming phase 3 \\
\code{co\_accessed} & Read together in the same step & Runtime tracking \\
\code{references} & Citation between nodes & Ingestion \\
\code{precedes, contradicts, supersedes, applies\_to} & Reserved for future use & --- \\
\bottomrule
\end{tabularx}
\end{table}

\subsection{Schema Property: One Trio Generalizes All Domains}

Because the three tables are disambiguated only by enums, the same
\code{CortexRouter.read}, \code{write}, \code{navigate}, vector-search SQL,
and graph-traversal CTEs work uniformly across Knowledge, Episodic,
Experience, Intelligence, and Runtime trees. New domains can be added by
extending the enum without altering the storage layer.

% =============================================================================
\section{The CORTEX Operations}
\label{sec:ops}

The \code{CortexRouter} class implements seven primitive operations. These
are surfaced to the agent both as planner step types (\op{Navigate},
\op{Read}, \op{Write}, \op{Recurse}, \op{Await\_Children}) and as direct
method calls (\op{Checkpoint}, \op{Assemble}).

\subsection{The Seven Primitives}

Table~\ref{tab:ops} lists the operations.

\begin{table}[h]
\centering
\caption{The seven primitive \textsc{Cortex} operations.}
\label{tab:ops}
\small
\begin{tabularx}{\linewidth}{lXX}
\toprule
\textbf{Op} & \textbf{Signature} & \textbf{Behavior}\\
\midrule
\op{Navigate} & \code{navigate(node\_id) $\to$ Viewport} & Move cursor; return current + children + parent + breadcrumb \\
\op{Read} & \code{read(node\_id, page=0) $\to$ NodeContent} & Return one page; promote node \code{pending} $\to$ \code{active} \\
\op{Write} & \code{write(parent\_id, type, title, content, summary, ...) $\to$ UUID} & Create child node; enforce invariants 1 and 2 \\
\op{Recurse} & \code{recurse(node\_id, task, result\_slot, ...) $\to$ (task\_node\_id, child\_run\_id)} & Create scoped-subtree child execution \\
\op{Await\_Children} & \code{await\_children(parent\_id) $\to$ \{slot: NodeSummary\}} & Collect results from completed child task nodes \\
\op{Checkpoint} & \code{checkpoint(tree\_id, progress, key\_facts, next\_steps) $\to$ UUID} & Compress context into a checkpoint node \\
\op{Assemble} & \code{assemble\_output(tree\_id, coherence\_pass=True) $\to$ str} & DFS the Output subtree; LLM-generate transition paragraphs \\
\bottomrule
\end{tabularx}
\end{table}

\subsection{The Viewport}

A viewport is the bounded representation of the agent's current position.
Token cost is bounded by $\text{max\_children} \times {\sim}40$ tokens per
child summary, plus the current node's summary, the parent's summary, and
the breadcrumb. The viewport is serialized as structured text:

\begin{lstlisting}[basicstyle=\ttfamily\scriptsize]
## Navigation Path
<breadcrumb>

## Current Node: <title>
Type / Status / Depth / Summary

## Children
  [1] <title> (<type>, <status>) -- <summary>
  ...

## Available CORTEX Operations
NAVIGATE(...) | READ(...) | WRITE(...) | RECURSE(...) |
AWAIT_CHILDREN() | CHECKPOINT(...)
\end{lstlisting}

\subsection{Bounded Context Property}
\label{sec:bounded}

\begin{property}[Bounded Viewport]
For any tree $T$ with $\text{max\_children} = M$, the prompt cost of
$\text{navigate}(n)$ for any node $n \in T$ is $O(M + \text{depth}(n))$
tokens, independent of $|T|$.
\end{property}

\begin{proof}[Justification]
The viewport contains the current node's summary ($O(1)$), at most $M$
child summaries ($O(M)$), one parent summary ($O(1)$), and a breadcrumb
whose length equals $\text{depth}(n)$.
\end{proof}

This property is the central design invariant: the agent's prompt cost is
decoupled from the tree's size. A tree with 100{,}000 nodes presents the
same per-step prompt cost as a tree with 10 nodes.

\subsection{Subtree Isolation}
\label{sec:isolation}

When a \code{CortexRouter} is constructed with
\code{scoped\_subtree\_root\_id}, every \code{\_get\_node} call validates
that the requested node is a descendant of the scope root via a single
recursive CTE:

\begin{lstlisting}[language=SQL]
WITH RECURSIVE ancestors AS (
  SELECT id, parent_id FROM cortex_nodes WHERE id = :node_id
  UNION ALL
  SELECT cn.id, cn.parent_id
  FROM cortex_nodes cn JOIN ancestors a ON cn.id = a.parent_id
)
SELECT 1 FROM ancestors WHERE id = :ancestor_id LIMIT 1;
\end{lstlisting}

Out-of-scope access raises an error.

\begin{property}[Subtree Isolation]
A child execution constructed with \code{scoped\_subtree\_root\_id}~$=s$
cannot read or write any node $n$ such that $s$ is not an ancestor of
$n$.
\end{property}

\subsection{Re-Clustering on Viewport Overflow}

When \op{Write} would cause a node's child count to exceed
\code{max\_children}, an asynchronous re-clustering routine moves the
first half of children under a new \code{group} node. The group inherits
the type of the moved children, preserving semantic homogeneity within
groups.

\subsection{Resumability}

A tree maintains \code{resume\_cursor\_id} updated by every \op{Navigate},
\op{Read}, and \op{Write}. On \code{resume\_tree(tree\_id)}:

\begin{enumerate}
  \item Status is flipped from \code{suspended} to \code{active}.
  \item The viewport at \code{resume\_cursor\_id} is loaded.
  \item The most recent \code{checkpoint} node under the cursor is
  returned.
  \item The execution engine's \code{\_\_completed\_steps\_\_} set is
  restored, causing already-finished steps to be skipped.
\end{enumerate}

This makes resumption deterministic and idempotent at step granularity.

% =============================================================================
\section{The Dreaming Consolidation Pipeline}
\label{sec:dreaming}

The \code{DreamingEngine} runs the cognitive sequence \textit{episode $\to$
observation $\to$ pattern $\to$ rule} as a three-phase background job.

\subsection{Trigger}

The pipeline runs when:
\begin{itemize}
  \item \code{dreaming\_worker} is enqueued (post-run hook or scheduled),
  and
  \item $\text{now}() - \text{experience\_tree.last\_consolidated\_at}
  \geq \text{CONSOLIDATION\_INTERVAL\_HOURS}$ (default 24), or
  \item \code{force=True}.
\end{itemize}

\subsection{Phase 1 --- Observation Extraction}

Inputs: episodes since \code{last\_consolidated\_at} (capped at
$\text{BATCH\_SIZE} = 20$), minimum
$\text{MIN\_EPISODES\_FOR\_DREAMING} = 5$.

The engine builds episode summaries
\code{\{id, task, status, tools\_used, cost\_usd, execution\_time\_ms\}}
and calls the LLM with \code{OBSERVATION\_EXTRACTION\_PROMPT}, asking for
observations across five categories: tool patterns, success factors,
failure patterns, cost patterns, and time patterns. Observations with
$\text{confidence} < 0.5$ are filtered. Survivors are written as
\code{OBSERVATION} nodes under the Experience tree's Observations section,
embedded via \code{EmbeddingService.embed\_node}, with
$\text{importance\_score} = \text{confidence}$.

\subsection{Phase 2 --- Pattern Recognition}

Observations are clustered by greedy embedding-cosine similarity
($> 0.75$). For each cluster of size $\geq 2$:

\begin{enumerate}
  \item The LLM is called with \code{PATTERN\_RECOGNITION\_PROMPT} over
  the cluster's summaries.
  \item A \code{PATTERN} node is written with
  \code{metadata\_extra = \{source\_observations, pattern\_strength,
  recurrence\_count, success\_correlation\}}.
  \item For each source observation, a
  \code{CortexEdge(edge\_type = "derived\_from",
  weight = 1/|cluster|,
  created\_by = "dreaming\_engine")} is created from pattern
  $\to$ observation.
\end{enumerate}

\subsection{Phase 3 --- Intelligence Distillation}

Strong patterns ($\text{strength} \geq 0.7$,
$\text{recurrence} \geq \text{MIN\_PATTERNS\_FOR\_DISTILLATION} = 2$)
are fed to the LLM with \code{INTELLIGENCE\_DISTILLATION\_PROMPT}
alongside existing rule summaries (for de-duplication). Output is an
array of typed rules: instruction, strategy, or preference. Each rule
becomes an \code{INSTRUCTION}, \code{STRATEGY}, or \code{PREFERENCE} node
under the appropriate Intelligence section. The tree's
\code{consolidation\_generation} is incremented and
\code{last\_consolidated\_at} is set.

\subsection{Properties of the Consolidation Pipeline}

\begin{itemize}
  \item \textbf{Monotonic accumulation.} Existing rules are not deleted;
  new generations add on top. Conflict resolution between generations is
  reserved for the \code{contradicts} and \code{supersedes} edge types
  (future work).

  \item \textbf{Confidence-weighted retrieval.} Rule retrieval at runtime
  (\code{IntelligenceTreeService.get\_applicable\_rules}) ranks by
  $\text{confidence} \times \text{cosine\_similarity}$, then bulk-updates
  \code{access\_count} and \code{last\_accessed\_at}. This implements a
  form of frequency-weighted importance.

  \item \textbf{Domain coupling via edges.} The \code{derived\_from} and
  (future) \code{generalizes} edges create a back-pointer chain from
  rule $\to$ pattern $\to$ observation $\to$ episode, making every rule
  fully traceable to the runs that produced it.
\end{itemize}

% =============================================================================
\section{The Semantic Graph Layer}
\label{sec:graph}

The \code{cortex\_edges} overlay provides associative cross-tree retrieval.

\subsection{Hybrid Search}

\code{SemanticGraphService.semantic\_graph\_search} proceeds in three
steps:

\begin{enumerate}
  \item Embed the query with $\text{task\_type} = $ \code{RETRIEVAL\_QUERY}.
  \item \textbf{Seed.} Run pgvector cosine search restricted to
  \code{(ct.entity\_id = :entity\_id OR ct.scope\_level IN ('app','tenant'))}
  and optionally to specified \code{memory\_domain}s. Return top-$k$
  seeds.
  \item \textbf{Expand.} For each seed, run BFS via a recursive CTE with
  cycle protection (\code{NOT (target\_node\_id = ANY(path))}), depth
  $\leq$ \code{expansion\_depth}, traversing edges with $\text{weight}
  \geq \text{min\_weight}$. Multiply weights along paths.
  \item \textbf{Re-rank.} Combine:
  $\text{combined\_score} = 0.7 \cdot \text{similarity}
  + 0.3 \cdot \text{edge\_weight}$. Sort descending, truncate to
  $2k$.
\end{enumerate}

This couples the precision of embedding search with the recall of graph
expansion.

\subsection{Auto-Edge Creation}

After any node is embedded, \code{create\_similarity\_edges(node\_id)}
finds up to 5 nodes with cosine similarity $\geq 0.85$ and creates
\code{semantic\_similar} edges weighted by the similarity score.

\subsection{Co-Access Tracking}

\code{track\_co\_access(node\_ids, run\_id)} creates pairwise
\code{co\_accessed} edges (initial weight 0.3) between all nodes the
agent reads in the same step. Subsequent re-traversals boost the weight
by $\text{BOOST\_ON\_TRAVERSAL} = 0.05$, capped at 1.0. Over time the
graph organizes along trajectories of nodes that are useful together.

\subsection{Maintenance}

A daily \code{graph\_maintenance\_worker}:

\begin{lstlisting}[language=SQL]
-- Decay weights of edges not traversed in 30 days
UPDATE cortex_edges
   SET weight = GREATEST(MIN_WEIGHT, weight * 0.95)
 WHERE last_traversed_at < NOW() - INTERVAL '30 days'
    OR last_traversed_at IS NULL;

-- Prune edges below minimum
DELETE FROM cortex_edges WHERE weight < MIN_WEIGHT;
\end{lstlisting}

Decay plus prune produces a self-trimming graph whose size scales with
active relevance rather than accumulated history.

% =============================================================================
\section{Scope Hierarchy}
\label{sec:scope}

\textsc{Cortex} implements a six-level scope hierarchy via the
\code{scope\_level} enum, allowing the same schema to represent
platform-wide, multi-tenant, per-user, per-agent, and per-run state
(Table~\ref{tab:scope}).

\begin{table}[h]
\centering
\caption{The six scope levels and their typical uses.}
\label{tab:scope}
\small
\begin{tabularx}{\linewidth}{lll>{\raggedright\arraybackslash}X}
\toprule
\textbf{Level} & \textbf{Value} & \textbf{Inherits to} & \textbf{Typical use}\\
\midrule
L0 & \code{app} & All tenants & Platform-wide instructions shared across all tenants \\
L1 & \code{partner} & Partner's tenants & Shared across tenants under a partner organization \\
L2 & \code{tenant} & Tenant's users & Tenant-level knowledge (company-wide playbook) \\
L3 & \code{user} & User's entities & Per-user preferences \\
L4 & \code{entity} & Entity's runs & Per-agent persistent trees \\
L5 & \code{runtime} & --- & One per execution (working memory) \\
\bottomrule
\end{tabularx}
\end{table}

The \code{semantic\_graph\_search} SQL clause
\code{(ct.entity\_id = :entity\_id OR ct.scope\_level IN ('app','tenant'))}
is the inheritance primitive: an entity sees its own trees plus shared
app/tenant trees, without duplicating data.

% =============================================================================
\section{Implementation}
\label{sec:impl}

\subsection{Stack}

\begin{itemize}
  \item \textbf{Language.} Python 3.10+.
  \item \textbf{Async ORM.} SQLAlchemy 2.x with asyncio support.
  \item \textbf{Database.} PostgreSQL with the \code{pgvector} extension.
  \item \textbf{Embeddings.} Pluggable; reference implementation uses
  Vertex AI (\code{text-embedding-005}, 768-dim) but the resolution path
  supports per-tenant override via an admin configuration table.
  \item \textbf{LLM.} Pluggable via an internal \code{LLMRouter}
  (reference implementation supports OpenAI, Anthropic, and Google
  providers).
  \item \textbf{Background jobs.} \code{arq} (Redis-backed) for
  \code{dreaming\_worker}, \code{graph\_maintenance\_worker},
  \code{cortex\_resume\_scheduled} (cron, every 5 minutes),
  \code{resume\_execution}, and \code{process\_document}.
\end{itemize}

\subsection{Execution Engine Integration}

The reference \code{ExecutionEngine} invokes \textsc{Cortex} in five
phases (C1--C5):

\begin{lstlisting}[basicstyle=\ttfamily\scriptsize]
C1: Create or resume CORTEX tree
C2: Memory assembly (assemble_memory facade routes v1/v2)
C3: Build context from viewport, knowledge subtree, context sources
C4: Locate working memory root
C5: Execute plan steps with per-step write_step, refresh_viewport,
    and periodic write_checkpoint
\end{lstlisting}

CORTEX-native step types (\op{Navigate}, \op{Read}, \op{Write},
\op{Recurse}, \op{Await\_Children}) are handled by
\code{CortexBridge.execute\_cortex\_step}. Conventional step types
(\op{Thought}, \op{Action}, \op{Tool\_Call},
\op{Child\_Entity\_Invocation}) are handled by the step executor, which
post-processes by writing each step result as a \code{finding} node and,
after tool calls, ingesting tool output as \code{knowledge} nodes with
provenance metadata.

\subsection{Performance Engineering}

Several optimizations were applied during development:

\begin{itemize}
  \item \textbf{Recursive CTEs replace iterative parent walks.}
  \code{\_build\_breadcrumb} and \code{\_is\_descendant\_of} use single
  recursive CTE queries, replacing $O(\text{depth})$ sequential SELECTs
  with one round-trip.

  \item \textbf{Incremental context-size tracking.}
  \code{update\_context\_size(key, old, new)} mutates an $O(1)$ counter
  on every \code{store\_step\_output}, replacing $O(n)$ full-context
  scans before auto-compaction checks.

  \item \textbf{Redis viewport caching.} \code{refresh\_viewport} caches
  the rendered viewport text under
  \code{cortex:viewport:\{tree.id\}:\{cursor\_id\}} with a 30~s TTL,
  avoiding redundant tree traversals when the cursor has not moved
  between steps.

  \item \textbf{Batch embedding.} \code{EmbeddingService.embed\_batch}
  batches at $\text{BATCH\_SIZE} = 100$ (the Vertex AI limit) with
  per-text truncation to 8000 characters.

  \item \textbf{UUID caching for long sessions.}
  \code{working\_root\_id}, \code{\_tree\_output\_root\_id}, and
  \code{\_tree\_id} are captured as locals before the first session
  commit to avoid \code{MissingGreenlet} errors from ORM-attribute
  expiry during long step loops.
\end{itemize}

\subsection{Safety and Multi-Tenancy}

\begin{itemize}
  \item \textbf{Tenant isolation.} Every \code{CortexRouter} is
  constructed with a \code{company\_id}; all SQL filters by it.
  \item \textbf{Subtree isolation.} See \S\ref{sec:isolation}.
  \item \textbf{Sensitive-context redaction.} A
  \code{\_SENSITIVE\_CONTEXT\_KEYS} set redacts substrings like
  \code{api\_key}, \code{secret}, \code{token}, \code{password} from
  \code{context\_state} before any database write.
  \item \textbf{Parent $\to$ child memory leak prevention.} When the
  execution engine spawns a child via \op{Child\_Entity\_Invocation},
  it strips parent-scoped memory keys (\code{\_\_memory\_\_},
  \code{\_\_episodic\_memory\_\_}, \code{\_\_semantic\_context\_\_},
  etc.) before child invocation.
  \item \textbf{Null-byte sanitization.} Auto-ingested context sources
  are stripped of \code{\textbackslash x00} bytes that PostgreSQL UTF-8
  columns reject.
  \item \textbf{REST authorization.} All \code{/api/v1/cortex/*}
  endpoints require authentication; \code{company\_id} is taken from
  the authenticated user, never from the request body.
\end{itemize}

% =============================================================================
\section{Discussion: Design Properties and Trade-Offs}
\label{sec:discussion}

\subsection{Properties}

Table~\ref{tab:properties} summarizes the system's design properties and
the mechanisms that produce them.

\begin{table}[h]
\centering
\caption{Design properties and the mechanisms that produce them.}
\label{tab:properties}
\small
\begin{tabularx}{\linewidth}{lX}
\toprule
\textbf{Property} & \textbf{Mechanism}\\
\midrule
Bounded per-step prompt cost & Viewport contains at most $M + \text{depth}(n)$ summaries \\
Cross-run continuity & Four persistent entity-scoped trees \\
Deterministic resumability & \code{resume\_cursor\_id} + checkpoint nodes + \code{\_\_completed\_steps\_\_} \\
Sub-agent isolation & \code{scoped\_subtree\_root\_id} with recursive-CTE ancestry check \\
Traceable provenance & \code{source\_ref} + \code{metadata\_extra.provenance\_chain} per node \\
Self-organizing graph & Edge weight boost on traversal, decay on inactivity, prune on threshold \\
Domain extensibility & Adding a domain or scope level is an enum extension, not a schema change \\
\bottomrule
\end{tabularx}
\end{table}

\subsection{Trade-Offs}

\begin{itemize}
  \item \textbf{Postgres-only.} The recursive CTEs and pgvector
  dependencies preclude SQLite or MongoDB backends. This is by design ---
  the structural query patterns are the architectural moat.

  \item \textbf{Async-only API.} Practical for modern Python agents but
  excludes synchronous codebases.

  \item \textbf{Operational complexity.} Running \textsc{Cortex} in
  production requires PostgreSQL, pgvector, Redis, and an arq worker.
  This is higher operational overhead than a library like Mem0 that runs
  in-process.

  \item \textbf{LLM cost during consolidation.} Dreaming makes three LLM
  calls per cycle. At typical observation/pattern volumes, this is
  sub-dollar per entity per day but is not free.

  \item \textbf{Greedy clustering in dreaming.}
  \code{\_cluster\_observations} is $O(n^2)$. Degrades past a few hundred
  observations; a HNSW-backed clustering pass is planned
  (\S\ref{sec:future}).

  \item \textbf{Synchronous re-clustering.} Triggered inside
  \code{write()} rather than queued; can briefly block step execution
  under high write pressure.

  \item \textbf{No node-level ACLs.} The unit of access is the tree.
  Per-node ACLs are reserved for future work.
\end{itemize}

\subsection{When \textsc{Cortex} Is and Is Not the Right Choice}

\textsc{Cortex} is well-suited to:

\begin{itemize}
  \item Long-horizon agents (multi-step, multi-hour, multi-day tasks).
  \item Agents that must produce structured output documents.
  \item Multi-tenant SaaS platforms with per-customer learning loops.
  \item Multi-agent systems where sub-agent isolation is critical.
\end{itemize}

\textsc{Cortex} is overkill for:

\begin{itemize}
  \item Single-turn chatbots needing only conversation history.
  \item Stateless RAG over a fixed document corpus (PageIndex or vanilla
  vector RAG suffice).
  \item Agents whose memory needs are exhausted by Mem0-style fact
  extraction.
\end{itemize}

% =============================================================================
\section{Limitations}
\label{sec:limits}

We list limitations honestly to delineate the system's current scope:

\begin{enumerate}
  \item \textbf{No empirical benchmarks yet.} This paper is a system
  description. Standard agent-memory benchmarks (LoCoMo, LongMemEval,
  AgentBench) are in the planning phase. No numerical claims about
  retrieval accuracy, downstream task completion, or token cost
  reduction are made here.

  \item \textbf{Dreaming has no automatic post-run trigger.}
  \code{dreaming\_worker} must be explicitly enqueued; a hook chained
  from \code{MemoryRouter.write\_episodic} is planned.

  \item \textbf{Memory scope \code{RUN\_SCOPED} is currently identical
  to \code{FULL}.} The filter to restrict retrieval to the current
  run's tree is not yet implemented.

  \item \textbf{Tree status remains \code{active} after run completion.}
  Long-term archival sweeps are not yet designed.

  \item \textbf{Concurrent writes under the same parent} rely on
  $\text{MAX}(\text{sibling\_order}) + 1$ without \code{SELECT FOR
  UPDATE}; the race window is small but not formally prevented.

  \item \textbf{v1 and v2 retrieval pipelines coexist.} The
  \code{MemoryAssemblyService} (v2) is implemented but the default
  per-entity pipeline is v1. Migration is gated by a feature flag on
  the entity capability config.

  \item \textbf{Reserved edge types} (\code{generalizes},
  \code{contradicts}, \code{supersedes}, \code{applies\_to}) are
  defined in the vocabulary but not yet emitted by the system.

  \item \textbf{No formal quality gate on distilled rules.}
  Low-confidence LLM outputs can pollute the Intelligence tree; a
  validation pass is planned.
\end{enumerate}

% =============================================================================
\section{Future Work}
\label{sec:future}

\subsection{Empirical Evaluation}

Planned experiments:

\begin{itemize}
  \item \textbf{Long-horizon task coherence} (analogous to Letta's
  30-day continuous-run benchmark).
  \item \textbf{Cross-session fact recall} (LoCoMo / LongMemEval
  comparison vs Mem0, Zep, Letta).
  \item \textbf{Cost per task} (Dreaming amortization vs naive
  context-stuffing baselines).
  \item \textbf{Sub-agent isolation correctness} (formal verification
  of Property 2 via fuzz testing).
  \item \textbf{Retrieval ablations}: pure embedding vs embedding +
  graph-1-hop vs embedding + graph-2-hop.
\end{itemize}

\subsection{Algorithmic Improvements}

\begin{itemize}
  \item Replace greedy $O(n^2)$ observation clustering with HNSW-backed
  batched clustering.
  \item Implement automatic post-run dreaming triggers gated by run
  count and time since last consolidation.
  \item Add LLM-based quality gate before writing Intelligence rules.
  \item Implement \code{contradicts}/\code{supersedes} edge emission
  during distillation; resolve conflicts at retrieval time.
\end{itemize}

\subsection{Operational Improvements}

\begin{itemize}
  \item Node-level ACLs for fine-grained sharing.
  \item Tree archival sweeps (move \code{complete} and
  $\text{inactive} > N$ days trees to a cold store).
  \item Async (queued) re-clustering instead of synchronous in
  \code{write()}.
  \item Backend-agnostic abstraction (although the design will remain
  Postgres-first).
\end{itemize}

\subsection{Theoretical Work}

\begin{itemize}
  \item Formalize the relationship between viewport bounded-context
  and downstream task accuracy.
  \item Connect the Dreaming consolidation pipeline to formal models
  of memory consolidation in cognitive science (e.g., the two-stage
  hippocampal-neocortical model).
  \item Investigate whether Intelligence rules form a useful ``policy
  distillation'' target for fine-tuning smaller models.
\end{itemize}

\subsection{Ecosystem}

Open-source as a standalone \code{pip}-installable package, with backend
plugins for OpenAI / Anthropic / litellm LLM providers and Vertex /
OpenAI / sentence-transformers embedding providers (planned).

% =============================================================================
\section{Conclusion}
\label{sec:conclusion}

We presented \textsc{Cortex}, an agent-memory architecture in which the
agent's complete cognitive state is a typed, persistent, navigable,
writable tree, and the language model's context window is a bounded
viewport onto that tree. \textsc{Cortex} unifies four memory domains and
a transient runtime workspace on a single ORM model trio, with seven
primitive operations, a three-phase background consolidation pipeline,
a semantic graph overlay, and a six-level scope hierarchy. Each
individual element draws on prior work --- MemGPT's tiering, PageIndex's
hierarchical navigation, RLM's recursive decomposition, Park et al.'s
reflection trees, CoALA's memory taxonomy --- but the integration into
one coherent runtime with one schema and one set of primitives is, to
our knowledge, novel. We have characterized the system's design
properties, documented its limitations, and outlined a path to
empirical evaluation. The architecture is in production use; the source
is available for inspection and extension.

% =============================================================================
\section*{Acknowledgments}

We thank the broader community working on LLM agent memory --- including
the teams behind MemGPT/Letta, Mem0, Zep, Cognee, Generative Agents,
PageIndex, and the CoALA framework --- for the foundational ideas this
work builds upon. Any errors or unsupported claims remain our own.

% =============================================================================
\begin{thebibliography}{99}

\bibitem{packer2023memgpt}
C.~Packer, S.~Wooders, K.~Lin, et~al.
\newblock {MemGPT: Towards LLMs as Operating Systems}.
\newblock \textit{arXiv preprint arXiv:2310.08560}, 2023.

\bibitem{bansal2025mem0}
P.~Bansal et~al.
\newblock {Mem0: Building Production-Ready AI Agents with Scalable
Long-Term Memory}.
\newblock \textit{arXiv preprint arXiv:2504.19413}, 2025.

\bibitem{park2023generative}
J.~S. Park, J.~C. O'Brien, C.~J. Cai, M.~R. Morris, P.~Liang, and
M.~S.~Bernstein.
\newblock {Generative Agents: Interactive Simulacra of Human Behavior}.
\newblock In \textit{Proceedings of UIST}, 2023.

\bibitem{sumers2023coala}
T.~R. Sumers, S.~Yao, K.~Narasimhan, and T.~L.~Griffiths.
\newblock {Cognitive Architectures for Language Agents}.
\newblock \textit{arXiv preprint arXiv:2309.02427}, 2023.

\bibitem{laird2022soar}
J.~E.~Laird.
\newblock {Introduction to the Soar Cognitive Architecture}.
\newblock \textit{arXiv preprint arXiv:2205.03854}, 2022.

\bibitem{vectifyai2025pageindex}
Vectify AI.
\newblock {PageIndex: Document Index for Vectorless, Reasoning-Based RAG}.
\newblock Technical report and open-source release, 2025.
\newblock \url{https://github.com/VectifyAI/PageIndex}.

\bibitem{zhang2025rlm}
A.~L. Zhang.
\newblock {Recursive Language Models}.
\newblock \textit{arXiv preprint arXiv:2512.24601}, 2025.

\bibitem{wang2025gmemory}
J.~Wang et~al.
\newblock {G-Memory: Tracing Hierarchical Memory for Multi-Agent Systems}.
\newblock In \textit{Advances in Neural Information Processing Systems},
2025.

\bibitem{byterover2026}
ByteRover Team.
\newblock {ByteRover: Agent-Native Memory Through LLM-Curated Hierarchical
Context}.
\newblock \textit{arXiv preprint arXiv:2604.01599}, 2026.

\bibitem{liu2025cogworkspace}
S.~Liu et~al.
\newblock {Cognitive Workspace: Active Memory Management for LLMs --- An
Empirical Study of Functional Infinite Context}.
\newblock \textit{arXiv preprint arXiv:2508.13171}, 2025.

\bibitem{lettafs2025}
Letta Inc.
\newblock {Benchmarking AI Agent Memory: Is a Filesystem All You Need?}
\newblock Letta blog, 2025.
\newblock \url{https://www.letta.com/blog/benchmarking-ai-agent-memory}.

\bibitem{memorysurvey2026}
{Anonymous}.
\newblock {Rethinking Memory Mechanisms of Foundation Agents in the
Second Half: A Survey}.
\newblock \textit{arXiv preprint arXiv:2602.06052}, 2026.

\end{thebibliography}

% =============================================================================
\appendix

\section{Glossary}
\label{app:glossary}

\begin{table}[h]
\centering
\small
\begin{tabularx}{\linewidth}{lX}
\toprule
\textbf{Term} & \textbf{Definition}\\
\midrule
\textsc{Cortex} & Cognitive Orchestrated Recursive Tree EXecution \\
Viewport & The bounded representation of the agent's current tree position \\
Cursor & \code{tree.resume\_cursor\_id} --- the node the agent is ``at'' \\
Domain & One of \{knowledge, episodic, experience, intelligence\} \\
Scope & Six-level hierarchy \{app, partner, tenant, user, entity, runtime\} \\
Dreaming & The three-phase background consolidation pipeline \\
Subtree anchor & A direct child of the runtime tree's root (Knowledge / Working / Output) \\
Working memory & The subtree where step \code{finding} nodes accumulate \\
Episode & A row in the persistent Episodic Tree representing one completed run \\
Observation & An LLM-extracted insight about a cluster of episodes \\
Pattern & A recurring behavior synthesized from clustered observations \\
Rule & An actionable directive distilled from strong patterns \\
\bottomrule
\end{tabularx}
\end{table}

\section{The Seven Operations --- Pseudocode Reference}
\label{app:pseudocode}

\begin{lstlisting}[language=Python]
# 1. NAVIGATE -- bounded read of one level of the tree
async def navigate(node_id) -> Viewport:
    node = await _get_node(node_id)
    tree.resume_cursor_id = node_id
    children = SELECT * FROM cortex_nodes
               WHERE parent_id = node_id ORDER BY sibling_order
    parent = await _get_node(node.parent_id) if node.parent_id else None
    breadcrumb = RECURSIVE_CTE_walk_to_root(node_id)
    return Viewport(node, children, parent, breadcrumb)

# 2. READ -- paged content access
async def read(node_id, page=0) -> NodeContent:
    node = await _get_node(node_id)
    if node.status == PENDING: node.status = ACTIVE
    tree.resume_cursor_id = node_id
    page_chars = tree.page_size_tokens * CHARS_PER_TOKEN
    return node.content[page * page_chars : (page + 1) * page_chars]

# 3. WRITE -- create child node; enforce invariants
async def write(parent_id, node_type, title, content, summary, ...) -> UUID:
    parent = await _get_node(parent_id)
    assert parent.summary, "Invariant 1: parent must have summary"
    if child_count(parent_id) >= tree.max_children:
        await _schedule_reclustering(parent_id)
    sibling_order = SELECT COALESCE(MAX(sibling_order), -1) + 1
                    FROM cortex_nodes WHERE parent_id = parent_id
    INSERT new node
    tree.total_nodes += 1
    tree.resume_cursor_id = new_node.id
    return new_node.id

# 4. RECURSE -- scoped subtree execution
async def recurse(node_id, task, result_slot) -> (UUID, UUID):
    task_node_id = await write(
        parent_id=node_id, node_type="task",
        content={task, result_slot, scoped_to: node_id}, ...
    )
    child_run = ExecutionRun(
        parent=current,
        input={cortex_tree_id, subtree_root_id: node_id, task, ...}
    )
    enqueue_async("execute_run", child_run.id)
    return (task_node_id, child_run.id)

# 5. AWAIT_CHILDREN -- collect completed child tasks
async def await_children(parent_node_id) -> Dict[str, NodeSummary]:
    children = SELECT * FROM cortex_nodes
               WHERE parent_id = parent_node_id
                 AND node_type = 'task' AND status = 'complete'
    return {c.metadata.result_slot: NodeSummary(c) for c in children}

# 6. CHECKPOINT -- compress context into a node
async def checkpoint(tree_id, progress, key_facts, next_steps) -> UUID:
    cursor = tree.resume_cursor_id
    return await write(
        parent_id=cursor, node_type="checkpoint",
        content={progress, key_facts, next_steps,
                 nodes_written, time_elapsed}, ...
    )

# 7. ASSEMBLE -- DFS output subtree + bridge paragraphs
async def assemble_output(tree_id, coherence_pass=True) -> str:
    sections = DFS_collect(tree.output_root_id, where status='complete')
    if coherence_pass and len(sections) > 1:
        bridges = LLM(
            prompt="Write transitions between sections", sections
        )
        return interleave(sections, bridges)
    return "\n\n".join(sections)
\end{lstlisting}

\section{Reproducibility}
\label{app:repro}

A reference implementation is available at the project repository (URL
to be released with the package) under an OSI-approved license. Postgres
schema migrations are shipped with the package. To reproduce the system
locally:

\begin{lstlisting}[basicstyle=\ttfamily\scriptsize]
# Prerequisites: PostgreSQL 14+ with pgvector, Redis
pip install cortex-memory
createdb cortex_dev && psql cortex_dev -c "CREATE EXTENSION vector;"
cortex-memory migrate --database-url postgresql+asyncpg://localhost/cortex_dev

# Minimal usage
from cortex_memory import CortexRouter, OpenAIEmbeddings, OpenAILLM
router = CortexRouter(
    db=session, tenant_id=..., embeddings=..., llm=...
)
tree = await router.create_tree(owner_id=..., task="Draft Q3 report")
node_id = await router.write(
    parent_id=tree.root_node_id, node_type="finding",
    title=..., summary=..., content=...
)
\end{lstlisting}

\end{document}
